\section{Two-step defects and exact quadratic ranks}\label{sec:quadratic}
Let $K$ be any field, let $A_v$ be finite-dimensional $K$-vector spaces,
and put $A=\prod_v A_v$ and $N=\sum_v\dim A_v$. Fix a finite simple
triangle-free subcubic graph $G$ and an operator $W:A\to A$ with blocks
$T_{vu}:A_u\to A_v$ vanishing off its edges. In particular $T_{vv}=0$.
Define
\[
 E=\operatorname{offdiag}(W^2),\qquad
 E_{ac}=\begin{cases}\sum_b T_{ab}T_{bc},&a\ne c,\\0,&a=c.\end{cases}
\]
All ranks below are ordinary $K$-linear ranks, including when $K$ has
characteristic two.

Call $(a,c)$ touched if $a\ne c$ and a two-edge path $a-b-c$ contains
an edge of $F=E(G)\setminus E(G_\square)$. Such a path has a unique
middle vertex; a second one would put both edges in a square. Define
the source-separated constraint map
\[
 \Phi:A\longrightarrow\prod_{c\in V}A,\qquad
 (\Phi x)_{c,a}=
 \begin{cases}E_{ac}x_c,&(a,c)\text{ touched},\\0,&\text{otherwise}.
 \end{cases}
\]
Unlike $Ex$, this map does not sum contributions from different source
coordinates. Set $C=\ker(I-W)$, $H=C\cap\ker\Phi$, and
$q=\rank(\Phi|_C)=\dim C-\dim H$.

\begin{theorem}[S: paid two-step stability]\label{thm:stability}
With the preceding definitions,
\begin{align*}
 2\dim H&\leq\sum_v\dim\pi_v(H),\\
 N&\leq2\rank(I-W)+2q,\qquad q\leq\rank\Phi\leq6\rank E,\\
 N&\leq2\rank(I-W)+12\rank E,\\
 N&\leq2\rank(I+W)+12\rank E.
\end{align*}
\formal{QuadraticBounds.lean}{subcubic_constrained_half,subcubic_paid_rank,subcubic_twoStep_stability}
\formal{TwoStepStability.lean}{constraintLoss_identity,constraintLoss_le}
\formal{TwoStepDefect.lean}{separatedConstraint_rank_le}
\end{theorem}
\begin{outline}
For an edge $uv\in F$, define $p_{uv}x$ to have coordinates
$T_{uv}x_v$ at $u$, $T_{vu}x_u$ at $v$, and zero elsewhere.
On the \emph{full} space $H$, the two-step constraints imply that this
vector again belongs to $H$, recovers across $uv$, and transmits zero
along every other incident edge. Thus $p_{uv}$ is an idempotent on $H$;
the idempotents for distinct unordered edges annihilate one another.
Each endpoint projection is injective on its edge factor, giving
exactly twice the factor dimension in total coordinate capacity.

This decomposition also splits capacity coordinate by coordinate. At
$u$ the operator $p_{uv}$ is induced by $T_{uv}T_{vu}$ on $\pi_u(H)$,
and similarly at $v$; at every other coordinate it is zero. Hence each
coordinate image is the direct sum of its edge-factor images and its
image from the residual $\ker\sum_{uv\in F}p_{uv}$. On that residual
all off-square transmissions vanish. Its actual coordinate observations
therefore recover on $G_\square$. Apply \cref{thm:square-good} to their
rank polymatroid, then add back the balanced edge factors. This proves
the intrinsic bound for $H$. No invariance is asserted for an arbitrary
subcode of $C$.

Rank--nullity gives $q+\dim H=\dim C$ and the paid bound. To compare
$\Phi$ with $E$, each block row of $E$ has at most six possible source
coordinates: three first neighbors, each with at most two nonreturning
continuations. Choose a basis from the \emph{actual} scalar row
functionals of $E$. Each chosen row has support in at most six blocks,
so the sum of the block-column ranks is at most $6\rank E$.
Projecting each column onto its touched targets cannot increase rank,
and source separation makes these ranks add. Therefore
$\rank\Phi\leq6\rank E$. Replacing $W$ by $-W$ leaves $E$ unchanged
and gives the final inequality.
\end{outline}

The actual-kernel restriction and coordinatewise splitting are checked in
\formal{EdgeProjections.lean}{edgeProjection_idempotent,edgeProjection_annihilate,edgeProjection_coordinatewise,edgeProjection_balanced}
For a full kernel of dimension $(1/2+\varepsilon)N$, the theorem requires
$q\geq\varepsilon N$ and $\rank E\geq\varepsilon N/6$. This rules out
a small two-step defect as the source of a fixed excess over half; it
does not prove that an arbitrary native recovery operator has small defect.

\begin{theorem}[Q: exact quadratic rank identity]\label{thm:quadratic-exact}
If $E=0$, equivalently $W^2$ is block diagonal, then
\[
 2\rank(I-W)=2\rank(I+W)=N+\rank(W^2-I).
\]
\formal{QuadraticBounds.lean}{subcubic_quadratic_exact}
\end{theorem}
\begin{outline}
First suppose $W^2=I$. The zero-defect case of \cref{thm:stability}
gives $\rank(I\pm W)\geq N/2$. The identity
$(I-W)(I+W)=0$ gives the reverse inequality for their sum. Both ranks,
and both kernel dimensions, are therefore exactly $N/2$, without an
eigenspace decomposition or a restriction on characteristic.

In general write $W^2=\operatorname{diag}(D_v)$ and set
$S_v=\ker(D_v-I)$. Then $S=\ker(W^2-I)=\prod_v S_v$.
Commutation of $W$ with $W^2$ shows that every actual block $T_{vu}$
maps $S_u$ into $S_v$. The restricted operator on $S$ is an involution
on the same graph support. Both $\ker(I-W)$ and $\ker(I+W)$ lie in $S$,
so their dimensions equal those of the restricted kernels, namely
$\dim S/2$. Finally $\dim S=N-\rank(W^2-I)$ and rank--nullity give
the displayed equalities.
\end{outline}

\begin{corollary}[Injective-subspace refinement and involutions]
\label{cor:quadratic-refinement}
Under $E=0$, if $W^2-I$ is injective on any subspace $U\leq A$, then
\[
 N+\dim U\leq2\rank(I-W),\qquad
 N+\dim U\leq2\rank(I+W).
\]
If $W^2=I$, both doubled ranks equal $N$.
\formal{QuadraticBounds.lean}{subcubic_quadratic_injective,subcubic_involution_ranks}
\end{corollary}
\begin{outline}
Injectivity gives $\dim U\leq\rank(W^2-I)$; apply
\cref{thm:quadratic-exact}. If $W^2=I$, that rank is zero.
\end{outline}

The injective-subspace statement includes any invariant summand on which
$W^2-I$ is invertible, but needs neither invariance nor a chosen Fitting
decomposition. The exact identity is stronger: it accounts for the
whole rank of $W^2-I$. For a construction, block-diagonal two-step
behavior cannot support a kernel above half ambient capacity, regardless
of widths, commutativity, or singular individual blocks. These are
additional sufficient hypotheses for the half bound; they do not imply
the original scalar matching bound for unrestricted vector codes.
